% ============================================================
% Ajustes de formato solicitados en la revisión del tutor.
%
% CÓMO USARLO: agregar la línea
%     % ============================================================
% Ajustes de formato solicitados en la revisión del tutor.
%
% CÓMO USARLO: agregar la línea
%     % ============================================================
% Ajustes de formato solicitados en la revisión del tutor.
%
% CÓMO USARLO: agregar la línea
%     % ============================================================
% Ajustes de formato solicitados en la revisión del tutor.
%
% CÓMO USARLO: agregar la línea
%     \input{preambulo_correcciones}
% en el preámbulo del documento principal (main.tex), después de
% cargar los paquetes habituales (caption, fancyvrb, float, etc.)
% y ANTES de \begin{document}.
% ============================================================

% ------------------------------------------------------------
% 1. Numeración de tablas y figuras por capítulo:
%    Tabla 2.1, Tabla 2.2, ..., Figura 3.1, Figura 3.2, ...
%    (el número de capítulo corresponde al contador de \section)
% ------------------------------------------------------------
\usepackage{chngcntr}
\counterwithin{table}{section}
\counterwithin{figure}{section}

% ------------------------------------------------------------
% 2. Tipo flotante "Código", numerado por capítulo (Código 2.1,
%    Código 3.1, ...). Reemplaza al antiguo "Listado".
%    Se usa en el texto con: \captionof{codigo}{Descripción}
%    Requiere el paquete caption (si el documento principal ya lo
%    carga, la línea \usepackage{caption} puede eliminarse).
% ------------------------------------------------------------
\usepackage{caption}
\DeclareCaptionType[within=section]{codigo}[Código][Índice de códigos]

% ------------------------------------------------------------
% 3. TikZ, usado por la figura del ciclo de desarrollo y por el
%    diagrama de clases del middleware (capítulo de Metodología).
% ------------------------------------------------------------
\usepackage{tikz}
\usetikzlibrary{arrows.meta, positioning, shapes.geometric, shapes.multipart, calc, fit}

% en el preámbulo del documento principal (main.tex), después de
% cargar los paquetes habituales (caption, fancyvrb, float, etc.)
% y ANTES de \begin{document}.
% ============================================================

% ------------------------------------------------------------
% 1. Numeración de tablas y figuras por capítulo:
%    Tabla 2.1, Tabla 2.2, ..., Figura 3.1, Figura 3.2, ...
%    (el número de capítulo corresponde al contador de \section)
% ------------------------------------------------------------
\usepackage{chngcntr}
\counterwithin{table}{section}
\counterwithin{figure}{section}

% ------------------------------------------------------------
% 2. Tipo flotante "Código", numerado por capítulo (Código 2.1,
%    Código 3.1, ...). Reemplaza al antiguo "Listado".
%    Se usa en el texto con: \captionof{codigo}{Descripción}
%    Requiere el paquete caption (si el documento principal ya lo
%    carga, la línea \usepackage{caption} puede eliminarse).
% ------------------------------------------------------------
\usepackage{caption}
\DeclareCaptionType[within=section]{codigo}[Código][Índice de códigos]

% ------------------------------------------------------------
% 3. TikZ, usado por la figura del ciclo de desarrollo y por el
%    diagrama de clases del middleware (capítulo de Metodología).
% ------------------------------------------------------------
\usepackage{tikz}
\usetikzlibrary{arrows.meta, positioning, shapes.geometric, shapes.multipart, calc, fit}

% en el preámbulo del documento principal (main.tex), después de
% cargar los paquetes habituales (caption, fancyvrb, float, etc.)
% y ANTES de \begin{document}.
% ============================================================

% ------------------------------------------------------------
% 1. Numeración de tablas y figuras por capítulo:
%    Tabla 2.1, Tabla 2.2, ..., Figura 3.1, Figura 3.2, ...
%    (el número de capítulo corresponde al contador de \section)
% ------------------------------------------------------------
\usepackage{chngcntr}
\counterwithin{table}{section}
\counterwithin{figure}{section}

% ------------------------------------------------------------
% 2. Tipo flotante "Código", numerado por capítulo (Código 2.1,
%    Código 3.1, ...). Reemplaza al antiguo "Listado".
%    Se usa en el texto con: \captionof{codigo}{Descripción}
%    Requiere el paquete caption (si el documento principal ya lo
%    carga, la línea \usepackage{caption} puede eliminarse).
% ------------------------------------------------------------
\usepackage{caption}
\DeclareCaptionType[within=section]{codigo}[Código][Índice de códigos]

% ------------------------------------------------------------
% 3. TikZ, usado por la figura del ciclo de desarrollo y por el
%    diagrama de clases del middleware (capítulo de Metodología).
% ------------------------------------------------------------
\usepackage{tikz}
\usetikzlibrary{arrows.meta, positioning, shapes.geometric, shapes.multipart, calc, fit}

% en el preámbulo del documento principal (main.tex), después de
% cargar los paquetes habituales (caption, fancyvrb, float, etc.)
% y ANTES de \begin{document}.
% ============================================================

% ------------------------------------------------------------
% 1. Numeración de tablas y figuras por capítulo:
%    Tabla 2.1, Tabla 2.2, ..., Figura 3.1, Figura 3.2, ...
%    (el número de capítulo corresponde al contador de \section)
% ------------------------------------------------------------
\usepackage{chngcntr}
\counterwithin{table}{section}
\counterwithin{figure}{section}

% ------------------------------------------------------------
% 2. Tipo flotante "Código", numerado por capítulo (Código 2.1,
%    Código 3.1, ...). Reemplaza al antiguo "Listado".
%    Se usa en el texto con: \captionof{codigo}{Descripción}
%    Requiere el paquete caption (si el documento principal ya lo
%    carga, la línea \usepackage{caption} puede eliminarse).
% ------------------------------------------------------------
\usepackage{caption}
\DeclareCaptionType[within=section]{codigo}[Código][Índice de códigos]

% ------------------------------------------------------------
% 3. TikZ, usado por la figura del ciclo de desarrollo y por el
%    diagrama de clases del middleware (capítulo de Metodología).
% ------------------------------------------------------------
\usepackage{tikz}
\usetikzlibrary{arrows.meta, positioning, shapes.geometric, shapes.multipart, calc, fit}
